% © 2026 Ing. David Jaroš — CC BY-NC-ND 4.0
%
% This work is licensed under a Creative Commons Attribution-NonCommercial-NoDerivatives
% 4.0 International License (CC BY-NC-ND 4.0).
%
% License History: Earlier drafts (up to v0.3) were released under CC BY 4.0.
% From v0.4 onward, all material is released under CC BY-NC-ND 4.0 to protect
% the integrity of the theoretical work during ongoing academic development.
%
% See LICENSE.md for full license text.

% UBT-AI-PROVENANCE-BEGIN
% schema: ubt-ai-provenance/v1
% tier: B_machine_verified
% ai_assistance: disclosed
% human_review: machine-verification
% editorial_responsibility: Ing. David Jaroš
% policy: ../../AI_PROVENANCE.md
% notice: Machine-verified against named sources or verifiers; individual attestation is not claimed.
% UBT-AI-PROVENANCE-END

\ifdefined\INCLUDEMODE
\else
  \documentclass[12pt,a4paper]{article}
  \usepackage[a4paper, margin=2.5cm]{geometry}
  \usepackage{amsmath, amssymb, amsthm}
  \usepackage{mathtools}
  \usepackage{hyperref}
  \usepackage{xcolor}
  \usepackage{booktabs}

  \newtheorem{theorem}{Theorem}[section]
  \newtheorem{lemma}[theorem]{Lemma}
  \theoremstyle{remark}
  \newtheorem{remark}[theorem]{Remark}

  \title{\textbf{B-Identification Proof Status}\\
         \large Bridge Between Vacuum-Polarization $B_0$ and $V_{\mathrm{eff}}(n)$ Coefficient $B$}
  \author{Ing. David Jaro\v{s}}
  \date{2026-05-10}

  \begin{document}
  \maketitle
\fi

\section{Goal}
We test whether the one-loop running coefficient $B_0$ from vacuum polarization
is rigorously identical to the coefficient $B$ in
\begin{equation}
V_{\mathrm{eff}}(n)=n^2-B\,n\ln n.
\end{equation}

\section{Starting point from canonical action}
From
\texttt{canonical/THEORY/canonical/canonical\_action.tex}, the kinetic sector is
\begin{equation}
\mathcal{L}_{\mathrm{kin}}=\mathrm{Tr}\!\left[(D_\mu\Theta)^\dagger(D^\mu\Theta)\right].
\end{equation}
With winding expansion along the compact phase direction,
\begin{equation}
\Theta(x,\psi)=\sum_{n\in\mathbb{Z}}\Theta_n(x)e^{in\psi},
\end{equation}
one obtains mode masses $m_n^2\sim n^2/R_\psi^2$ and one-loop determinants of
Coleman--Weinberg type:
\begin{equation}
\Delta W_{\mathrm{1loop}}\sim \frac{1}{2}\sum_n\int\frac{d^4p}{(2\pi)^4}\ln\left(p^2+m_n^2\right).
\end{equation}

\section{What is proved vs. not proved}
\begin{itemize}
  \item \textbf{Proved:} $B_0$ controls logarithmic running in
  $1/\alpha(\mu)$ through vacuum polarization.
  \item \textbf{Not proved:} that the same coefficient must appear as the
  multiplier of the $n\ln n$ term in $V_{\mathrm{eff}}(n)$.
\end{itemize}

Existing canonical/research derivations already indicate that direct
one-loop determinants naturally produce different $n$-dependence unless an
additional mechanism is supplied; therefore the equality $B=B_0$ is currently
an \textbf{open identification problem}, not a theorem.

\section{Lemma status}
\begin{lemma}[Current bridge status]
The identity
\begin{equation}
B\stackrel{?}{=}B_0
\end{equation}
is currently \textbf{not proved} from $S[\Theta]$ in the canonical alpha chain.
\end{lemma}

\begin{proof}
$B_0$ comes from the photon two-point function (running-coupling sector), while
$B$ in $V_{\mathrm{eff}}(n)$ comes from winding-mode effective action terms.
A derivation equating these coefficients requires an explicit theorem mapping
those two sectors. No such theorem currently exists in the audited canonical
files; therefore the identification remains open.
\end{proof}

\section{Consequence for current alpha program}
\begin{itemize}
  \item If one uses the audited minimal-scalar loop count
  ($N_{\mathrm{eff}}=3$ in \texttt{canonical/n\_eff/step2\_AUDIT.tex}), then
  $B_0=2\pi$ and $n_*\approx 10.54$ in the currently used stationarity equation.
  \item Therefore the route to $n_*=137$ is \textbf{conditional-weak} until both
  (i) multiplicity audit gaps and (ii) $B_0\to B$ bridge gap are closed.
\end{itemize}

\ifdefined\INCLUDEMODE
\else
  \end{document}
\fi
